% from ega4-9-fr.tex
\oldpage[IV-3]{70}

supposer $A$ intègre et $X$ réduit, donc intègre puisqu'il est irréductible. On
sait (4.1.2) qu'il existe un ouvert $W$ dense dans $X_\eta$ et un
$k(\eta)$-morphisme \emph{fini et surjectif}
$h:W\to\mathbf V((k(\eta))^n)=\operatorname{Spec}(B')$, où
$B'=k(\eta)[T_1,\ldots,T_n]$ ($T_i$ indéterminées). Si $V$ est un ouvert de
$X$ tel que $V\cap X_\eta=W$, on sait (9.5.3) que pour $s$ voisin de $\eta$
dans $S$, $V_s$ est un ouvert dense dans $X_s$, et l'on peut par suite (4.1.1.3)
se borner au cas où $V=X$, $W=X_\eta$. Posons
$B=A[T_1,\ldots,T_n]$ et $Y=\operatorname{Spec}(B)=\mathbf V_A^n$, de sorte
que $\operatorname{Spec}(B')=Y_\eta$~; il résulte de (8.8.2, (i)) et
(8.10.5, (vi) et (x)), appliqués suivant la méthode de (8.1.2, \emph{a)}),
qu'en remplaçant au besoin $S$ par un voisinage de $\eta$, on peut supposer que
$h=g_\eta$, où $g:X\to Y$ est un morphisme fini surjectif~; autrement dit, on
a $X=\operatorname{Spec}(C)$, où $C$ est une $B$-algèbre finie et
l'homomorphisme $B\to C$ correspondant à $g$ est \emph{injectif}~; comme $C$
est un anneau intègre, $C$ est donc un $B$-module de type fini \emph{sans
torsion}, et $C_\eta=C\otimes_A k(\eta)$ est donc un module de type fini sans
torsion sur $B_\eta=B'=B\otimes_A k(\eta)=k(\eta)[T_1,\ldots,T_n]$ (étant un
module de fractions dont les dénominateurs sont dans
$A-\{0\}\subset B-\{0\}$). Il résulte donc de (9.4.8) qu'il y a un voisinage
$U$ de $\eta$ dans $S$ tel que pour $s\in S$,
$C_s=C\otimes_A k(s)$ soit un module de type fini \emph{sans torsion} sur
$B_s=B\otimes_A k(s)=k(s)[T_1,\ldots,T_n]$, et en particulier l'homomorphisme
$g_s:B_s\to C_s$ est injectif. Comme aucun élément de $B_s$ n'est diviseur de
0 dans $C_s$, pour tout idéal premier minimal $\mathfrak p_i$ de $C_s$ (dont
les éléments sont diviseurs de 0 dans $C_s$), on a nécessairement
$\mathfrak p_i\cap B_s=0$, donc l'homomorphisme canonique
$B_s\to C_s/\mathfrak p_i$ est injectif. On en déduit que pour chaque
composante irréductible $Z_i=\operatorname{Spec}(C_s/\mathfrak p_i)$ de $X_s$,
la restriction à $Z_i$ de $g$ est un morphisme fini et dominant $Z_i\to Y_s$,
donc \emph{surjectif} (\textbf{II}, 6.1.10). On conclut par suite de (4.1.2)
que $\dim Z_i=n$, ce qui achève la démonstration.

\begin{remark}[9.5.7]
\label{IV.9.5.7-fr}
On aura soin de noter que sous les hypothèses de (9.5.6) il peut se faire que
$X_s$ \emph{ne soit irréductible pour aucun} $s\neq\eta$ dans un voisinage de
$\eta$, autrement dit, la propriété «~$X$ est un $k$-préschéma algébrique
irréductible~» \emph{n'est pas constructible}. Prenons par exemple
$S=\operatorname{Spec}(k[T])$, où $k$ est un corps algébriquement clos, $T$ une
indéterminée~; on a donc $k(\eta)=K=k(T)$. Soit $L$ une extension finie
séparable de $K$ de degré $>1$, et soit $X$ la \emph{fermeture intégrale} de
$S$ dans $L$ (\textbf{II}, 6.3.4)~; on a donc $X=\operatorname{Spec}(B)$, où
$B$ est la fermeture intégrale de $k[T]$ dans $L$. On sait que $B$ est un
anneau de Dedekind, et que tous les idéaux maximaux de $k[T]$, sauf un nombre
fini, sont non ramifiés dans $B$~; comme en outre le corps résiduel de tout
idéal maximal de $B$ est nécessairement $k$ (puisque c'est une extension finie
de $k$), on voit que pour presque tous les idéaux maximaux $\mathfrak j_s$ de
$k[T]$, $B_s=B/\mathfrak j_sB$ est composé direct de $[L:K]$ corps isomorphes
à $k$, autrement dit $X_s$ n'est pas irréductible, bien que
$X_\eta=\operatorname{Spec}(L)$ le soit. Le même exemple montre que la
propriété «~$X$ est un $k$-préschéma algébrique intègre~» n'est pas
constructible. Enfin il en est de même de la propriété «~$X$ est un
$k$-préschéma algébrique réduit~». Il suffit pour le voir de prendre encore
$S=\operatorname{Spec}(k[T])$, où cette fois $k$ est un corps algébriquement
clos de caractéristique $p>0$, et pour $X$ la fermeture intégrale de $S$ dans
$L=K^{1/p}$ (où $K=k(T)$), de sorte que $X=\operatorname{Spec}(B)$ avec
$B=k[T^{1/p}]$ ($k$ étant parfait)~; tout idéal maximal de $k[T]$ est de la
forme $(T-\alpha)$ avec $\alpha\in k$~; l'unique idéal de $B$ au-dessus de
l'idéal $(T-\alpha)$ est l'idéal principal
$(T^{1/p}-\alpha^{1/p})$ et il est immédiat

